\documentclass[11pt]{article}
\usepackage{preamble}
\setlength{\parskip}{4pt}

\begin{document}

\daybanner{AP Statistics \textperiodcentered\ Unit 5 \textperiodcentered\ Topic 5.5 \textperiodcentered\ TEACHER KEY}{Which Regression Line Fits?}

\sectionbanner{CED TETHER}

{\small
\begin{itemize}[leftmargin=1.5em,itemsep=2pt,topsep=2pt]
  \item \textbf{Skill 2.C} --- Calculate summary statistics, relative positions of points within a distribution, correlation, and predicted response.
  \item \textbf{Skill 4.B} --- Interpret statistical calculations and findings to assign meaning or assess a claim.
  \item \textbf{EU DAT-1} --- Regression models may allow us to predict responses to changes in an explanatory variable.
  \item \textbf{LO DAT-1.G} --- Estimate parameters for the least-squares regression line model. [Skill 2.C]
  \item \textbf{EK DAT-1.G.1} --- The least-squares regression model minimizes the sum of the squares of the residuals and contains the point ($\bar{x}$, $\bar{y}$).
  \item \textbf{EK DAT-1.G.2} --- The slope, b, of the regression line can be calculated as b = r(s\_y / s\_x), where r is the correlation between x and y, s\_y is the sample standard deviation of the response variable, and s\_x is the sample standard deviation of the explanatory variable.
  \item \textbf{EK DAT-1.G.3} --- Sometimes, the y-intercept of the line does not have a logical interpretation in context.
  \item \textbf{EK DAT-1.G.4} --- In simple linear regression, r\textsuperscript{2} is the square of the correlation, r. It is also called the coefficient of determination. r\textsuperscript{2} is the proportion of variation in the response variable that is explained by the explanatory variable in the model.
  \item \textbf{LO DAT-1.H} --- Interpret coefficients for the least-squares regression line model. [Skill 4.B]
  \item \textbf{EK DAT-1.H.1} --- The coefficients of the least-squares regression model are the estimated slope and y-intercept.
  \item \textbf{EK DAT-1.H.2} --- The slope is the amount that the predicted y-value changes for every unit increase in x.
  \item \textbf{EK DAT-1.H.3} --- The y-intercept value is the predicted value of the response variable when the explanatory variable is equal to 0. The formula for the y-intercept, a, is a = $\bar{y}$ - b$\bar{x}$.
\end{itemize}
}
\textbf{Essential question:} How do several regression properties distinguish lines that share the same mean point?\par
\textbf{DOK-3 skill:} 4.B \textperiodcentered\ \textbf{FRQ pattern:} {\small\texttt{shared-\allowbreak{}mean-\allowbreak{}point-\allowbreak{}which-\allowbreak{}line-\allowbreak{}is-\allowbreak{}lsrl}}\par
\textbf{DOK rationale:} Part (c) weighs competing claims using the day's numerical or graphical evidence and explains a limitation or consequence; the judgment requires linked reasoning beyond a calculation.\par

\textbf{Self-paced bonus sheet.} Students take it when they choose and turn it in when they finish. Everything they need is printed on the sheet. Score part (c) only, E / P / I, by hand --- never AI-graded or auto-scored. (a) and (b) are the ladder up; use them to see where a thin (c) came from.\par

\textbf{Questions to ask}
\begin{itemize}[leftmargin=1.5em,itemsep=2pt,topsep=2pt]
  \item Which number or visible feature supports your claim?
  \item Why does the other claim go beyond that evidence?
  \item What would you tell someone who chose the other claim?
\end{itemize}

\begin{callout}{\IconWarn\ LOOK FOR}{calloutred}
\begin{itemize}[leftmargin=1.5em,itemsep=2pt,topsep=2pt]
  \item A choice with copied numbers but no explanation of how they support it.
  \item Treating a model or average as a guaranteed individual outcome.
\end{itemize}
\end{callout}

\sectionbanner{SCORING PART (c) --- E / P / I}

\textbf{Expected elements}
\begin{itemize}[leftmargin=1.5em,itemsep=2pt,topsep=2pt]
  \item \textbf{reasoning-1} (required) --- Keeps Quinn's line using the correct SD ratio and slope
  \item \textbf{reasoning-2} (required) --- Explains why the shared mean point is insufficient to choose the LSRL
  \item \textbf{reasoning-3} (required) --- Keeps Rowan's 64 percent and interprets response variation in context
\end{itemize}

{\small\renewcommand{\arraystretch}{1.25}
\noindent\begin{tabular}{|p{0.5in}|p{5.9in}|}
\hline
\textbf{E} & Includes all three required reasoning elements above, with correct contextual evidence. \\ \hline
\textbf{P} & Includes at least one substantive correct reasoning link above, but omits or misstates another required element. \\ \hline
\textbf{I} & Includes no substantive correct reasoning link above; an unsupported choice or copied number alone is insufficient. \\ \hline
\end{tabular}}

\textbf{Common mistakes}
\begin{itemize}[leftmargin=1.5em,itemsep=2pt,topsep=2pt]
  \item Reversing the standard-deviation ratio
  \item Treating the mean point as sufficient
  \item Using $r$ as explained variation
\end{itemize}

\clearpage
\focusbanner{THE PROBLEM --- annotated}

Suppose a solar-farm analyst uses sunshine hours $x$ to predict daily generation $y$ in megawatt-hours (MWh). The table summarizes 12 observed days.\par\smallskip \textbf{Quinn:} ``Using $b=r(s_y/s_x)$ gives $\hat y=5.6+3.2x$. It passes through $(7,28)$, so it must be the LSRL. It explains 80 percent of the variation.''\par \textbf{Rowan:} ``Using $b=r(s_x/s_y)$ gives $\hat y=26.6+0.2x$. It also passes through $(7,28)$, so it is equally good. Explained variation is 64 percent.''\par
\begin{center}
\textbf{Solar-farm summary (12 observed days)}\par\smallskip
\resizebox{0.98\linewidth}{!}{%
\begin{tabular}{|c|c|c|c|c|c|}
\hline
$x$ range & $\bar x$ & $\bar y$ & $s_x$ & $s_y$ & $r$ \\ \hline
4--10 h & 7 h & 28 MWh & 1.5 h & 6 MWh & 0.80 \\ \hline
\end{tabular}}
\end{center}
\begin{firsttakebox}{FIRST TAKE --- one sentence before you work the parts}
The two lines agree at 7 sunshine hours. Does that make them equally useful at other times? Give your first thought in one sentence.\par
\answer{Not graded. Credit a committed clause-level comparison supported by one property.}
\end{firsttakebox}

\begin{callout}{\IconBook\ MAKE THE REGRESSION PROPERTIES AGREE (keep on the page)}{calloutgreen}
The least-squares regression line (LSRL) minimizes the sum of squared residuals and contains
$(\bar x,\bar y)$. Its slope is $b=r(s_y/s_x)$, its intercept is
$a=\bar y-b\bar x$, and $r^2$ is the proportion of response variation explained by the
linear model. Interpret slope as a predicted response change per unit of $x$. The intercept
predicts at $x=0$, which may be outside the observed range and unsupported in context.
\end{callout}

\dokbadge{1}~\textbf{(a)} State the mean point that every least-squares regression line must contain.\par
\answer{The mean point is $(\bar x,\bar y)=(7,28)$: 7 sunshine hours and 28 MWh.}

\dokbadge{2}~\textbf{(b)} Calculate the regression slope and intercept, and $r^2$. Give the slope's units.\par
\answer{$b=0.80(6/1.5)=3.2$ MWh per hour; $a=28-3.2(7)=5.6$ MWh; $r^2=(0.80)^2=0.64$.}

\dokbadge{3}~\textbf{(c)} Which parts of Quinn's and Rowan's claims should a report keep? Use the slope formula to choose a line, explain why their shared mean point does not make both lines correct, and interpret the correct percent of variation explained. Write 3--4 sentences.\par
\answer{Keep Quinn's line, $\hat y=5.6+3.2x$, because it uses the correct response-SD over explanatory-SD ratio. Both lines pass through $(7,28)$, but only Quinn's has the regression slope; passing through the mean point alone is not enough to minimize squared residuals. Keep Rowan's $64\%$: the linear model explains $64\%$ of the variation in daily generation. Quinn's $80\%$ confuses correlation with its square.}

\end{document}
